\section*{Abstract}

Freshwater transport by the overturning at the southern boundary of the
Atlantic, $\Fovs$, is the standard indicator of whether the salt-advection
feedback is self-reinforcing, and its recent decline in ocean reanalyses is
increasingly read as a warning that the overturning is weakening. We show that
it is not. Once the net throughflow is removed, $\Fovs$ factorises exactly as
$-T\,\Delta S/S_0$, with $T$ the northward-limb volume transport and $\Delta S$
the transport-weighted salinity difference between the limbs, both measured
directly from the section. At 34.5\degS{} $\Fovs$ changes at 36\% (ORAS5) and
52\% (GLORYS12V1) of its own magnitude per decade while $T$ changes at 2.6\%
and 0.6\%. The transport contributes 1\% to 18\% of the trend depending on how
the overturning is defined, in most configurations with the opposite sign, and
stays below 25\% in 97\% of all admissible epoch pairs in the two declining
products. What the decline
records instead is a widening limb salinity contrast, carried by the northward
limb; the EN4 observational analysis reproduces that asymmetry in a salinity field no
model integrates. The products disagree on the rest: ECCO-V4r4 has no trend to
attribute, and in ORAS5 the split between salinity and flow structure is not
resolved, with a third of the widening attributable to the flow.

\section*{Introduction}

The Atlantic meridional overturning circulation transports salt as well as
heat, and the resulting salt-advection feedback is the classical route to
multiple equilibria in this basin\cite{Stommel1961,Rahmstorf1996,Rahmstorf2002}.
The diagnostic used to test whether that feedback is self-reinforcing is
$\Fovs$, the freshwater transport carried by the overturning across the
southern boundary\cite{deVries2005,Huisman2010,Dijkstra2007}. When $\Fovs<0$
the overturning imports freshwater into the basin, so a weakening circulation
freshens the Atlantic and weakens it further; when $\Fovs>0$ the same
perturbation is damped. In box models and in coarse-resolution ocean models the
sign of $\Fovs$ separates a monostable from a bistable regime, and its value has
been used to argue that most climate models are biased towards the stable
side\cite{Weijer2019,Mecking2017,vanWestenDijkstra2024}.

That correspondence is not exact. Extended box models\cite{Cimatoribus2014}
and hysteresis experiments in a state-of-the-art coupled
model\cite{vanWestenDijkstra2023} show that the collapsed branch can persist
well outside the $\Fovs<0$ range, so the sign indicates where the freshwater
budget sits rather than proving that a second equilibrium exists. The
\emph{transient} behaviour of $\Fovs$ is a further step removed. $\Fovs$
reflects the salinity contrast between the limbs at the section, which responds
to North Atlantic anomalies only after they have been advected south, and an
AMOC collapse in a coupled model need not coincide with a minimum in
$\Fovs$\cite{vanWesten2025,Vanderborght2025}. Despite this, a declining $\Fovs$
is increasingly read as a physics-based early warning that the circulation is
weakening\cite{vanWesten2024,Zhu2020,Zhu2023}, and it is that inference we test.

The inference is not safe, for a reason visible in the definition. Write
$\Fovs$ as the product of an overturning transport and a salinity contrast.
Wherever $\Fovs<0$ the contrast is positive, so a weakening overturning
\emph{raises} $\Fovs$ rather than lowering it, and a falling $\Fovs$ cannot by
itself be evidence of a weakening circulation\cite{Vanderborght2025}. (The
argument applies only where $\Fovs<0$ throughout; ORAS5 has 7 years of 68 in
which it does not hold.) Whether the observed decline is carried by the
transport or by the salinity field is an empirical question, and it determines
how much of the recent literature on $\Fovs$ trends can be interpreted at all.

Answering it requires a factorisation whose two factors are separately
measurable. Defining the salinity contrast as $-S_0\Fovs/\Psi$ for some scalar
overturning strength $\Psi$ does not qualify: the second factor is then the
first divided out, it absorbs every change in the vertical structure of the
velocity as well as every change in salinity, and the resulting attribution is
fixed by algebra rather than by data. We therefore use the limb factorisation.
Removing the net section throughflow, which any transport-consistent estimate
of $\Fovs$ requires, forces the depth integral of the section velocity to
vanish, so the northward and southward limbs carry equal and opposite volume
transport $T$, and splitting the depth integral at the sign changes of the
velocity gives
\begin{equation}
\Fovs \;=\; -\frac{1}{S_0}\,T\,\Delta S,
\qquad \Delta S \;=\; S_{\uparrow}-S_{\downarrow},
\label{eq:factorisation}
\end{equation}
exactly and with no residual, where $S_{\uparrow}$ and $S_{\downarrow}$ are the
transport-weighted mean salinities of the northward and southward limbs.
Equation~\eqref{eq:factorisation} is not a linearisation and not a fit. Both
factors are computed directly from the section profile, $T$ is linear in the
velocity field and so carries none of the sampling bias of a streamfunction
maximum, and $\Delta S$ is a salinity difference that can be plotted, compared
between products, and in principle measured by hydrography.

We evaluate this at 34.5\degS{} in three ocean reanalyses recomputed here from
monthly fields, ORAS5, GLORYS12V1 and ECCO-V4r4, spanning two ocean models and
three assimilation methods, and we ask what the answer implies for the
interpretation of CMIP6 freshwater transport biases.
